% -------------------------------------------------
% VERSION TRACKER — No modificar este bloque
% Versión:    Tesis Humanizada
% Archivo:    Final/TrabajosFuturos.tex
% Última mod:  2026-05-08T08:45:26
% Git commit:  cc9c766f @ main
% Audiencia:   general
% Verificar antes de editar: bash check_tesis_version.sh overleaf_tesis_humanizado/Final/TrabajosFuturos.tex
% -------------------------------------------------
% Revisado contra el código el 26-sep-2026: se retiraron propuestas que ya están
% implementadas (despliegue en OCI, CI/CD, canales de notificación, órdenes de compra,
% cifrado de claves de API) y se acotaron las que existen solo en parte.


\section{Trabajos Futuros}

Las propuestas se agrupan en cinco frentes. El primero recoge lo que quedó pendiente del caso de estudio y es el más directamente útil para la empresa; los demás amplían el sistema en infraestructura, inteligencia artificial, funcionalidades de negocio, y experiencia de usuario, seguridad e investigación. En cada caso se indica qué existe hoy, para distinguir lo que falta de lo que ya se construyó.

\subsection{Requerimientos del caso de estudio pendientes}
Del levantamiento de requerimientos (Capítulo \ref{chap:requerimientos}) quedaron por completar:
\begin{itemize}
    \item \textbf{Desglose del consumo por destino (RF-15)}: registrar cuánto se usa en esponja, espolvoreo y cada línea de producto, que es la causa de las diferencias que la empresa detecta en el cierre.
    \item \textbf{Factura y ficha de calidad en el ingreso (RF-03)} y un \textbf{registro de no conformidades en la recepción (RF-14)}. Hoy el proveedor queda asociado a la orden de compra y un lote puede pasar a cuarentena, pero la factura, la ficha y la novedad no tienen un registro propio.
    \item \textbf{Lectura de los códigos de barras de fábrica} de los insumos, además de las etiquetas QR propias del sistema.
    \item \textbf{Conversión automática de presentaciones} (bultos y cajas) a la unidad base en kilogramos (RNF-01); hoy la conversión la hace el usuario.
\end{itemize}

\subsection{Infraestructura y despliegue}
\begin{itemize}
    \item \textbf{Copias de seguridad fuera del servidor}: el volcado nocturno de la base de datos y los respaldos manuales previos a cambios delicados (Capítulo \ref{chap:implementacion}) quedan en el mismo servidor, de modo que una falla de la instancia los perdería. Se propone enviarlos a otro almacenamiento de capa gratuita (por ejemplo, OCI Object Storage), incluir los archivos subidos y probar periódicamente la restauración.
    \item \textbf{Integración continua más estricta}: GitHub Actions ya revisa la sintaxis, ejecuta PHPUnit y despliega solo si las pruebas pasan; faltan la suite Playwright contra una instancia temporal y el análisis estático (PHPStan y ESLint) como condición para desplegar.
    \item \textbf{Contenedores Docker}: actualizar el \texttt{Dockerfile} y el \texttt{docker-compose.yml} del repositorio, que conservan valores de una etapa anterior, y validarlos en la integración continua.
    \item \textbf{Multi-tenencia}: aislar los datos de varias PYMEs en una sola instancia mediante un identificador de inquilino (\texttt{tenant\_id}).
    \item \textbf{Modelo local afinado}: el respaldo local \texttt{qwen3.5:9b} tarda de 30 a 60 segundos por respuesta (Capítulo \ref{chap:implementacion}); afinar un modelo más pequeño con consultas del dominio, o usar un servidor con GPU, lo haría más rápido.
\end{itemize}

\subsection{Inteligencia artificial y RAG}
\begin{itemize}
    \item \textbf{Consultas con herramientas}: que el modelo pida los datos que necesita mediante funciones de solo lectura sobre la base de datos, con el sistema como filtro de seguridad y una verificación automática de que cada cifra de la respuesta aparezca en los datos consultados.
    \item \textbf{Búsqueda semántica con embeddings}, para responder por significado (``¿qué grasas tenemos?'') sin conocer el nombre exacto de cada insumo.
    \item \textbf{Memoria conversacional en la recuperación}: el modelo recibe los dos últimos turnos, pero la clasificación usa solo la pregunta actual, así que ``¿y de azúcar?'' tras ``¿cuánta harina queda?'' se interpreta sin contexto. Se propone heredar la intención y el material del turno anterior.
    \item \textbf{Mejor pronóstico y detección de anomalías}: la proyección actual es un promedio móvil de cuatro semanas (Sección \ref{sec:imp_proyeccion}), que no anticipa temporadas ni tendencias; se propone probar modelos de series temporales con el historial del Kardex y señalar consumos atípicos, entradas duplicadas o ajustes frecuentes.
    \item \textbf{Evaluación continua del asistente}: programar la batería de \texttt{rag:evaluar} (Sección \ref{sec:pru_asistente}) y alertar cuando la tasa de acierto baje de un umbral.
\end{itemize}

\subsection{Funcionalidades de negocio}
\begin{itemize}
    \item \textbf{Cotizaciones} previas a las órdenes de compra, que hoy ya cubren proveedor, aprobación y recepción con lote, vencimiento y bodega.
    \item \textbf{Catálogo de categorías} propio de cada empresa, para renombrarlas y fusionarlas; hoy la categoría es un texto del insumo.
    \item \textbf{Integración con proveedores} (precios, disponibilidad y orden de reposición al cruzar el punto de reorden).
    \item \textbf{Módulo de producción (BOM)}: recetas que descuenten ingredientes por FEFO al producir y calculen el costo real; resolvería buena parte del RF-15.
    \item \textbf{Contabilidad integrada} y \textbf{facturación electrónica DIAN}.
\end{itemize}

\subsection{Experiencia de usuario, seguridad e investigación}
\begin{itemize}
    \item \textbf{Registro sin conexión}: la aplicación ya funciona en el celular y su \textit{Service Worker} guarda en caché los archivos estáticos; falta registrar movimientos sin red y sincronizarlos respetando FEFO y la inmutabilidad del Kardex.
    \item \textbf{Notificaciones \textit{push}} del navegador, además de las alertas en la aplicación, Telegram y correo; \textbf{identidad visual por empresa} conservando el contraste WCAG 2.1 AA de los seis temas; \textbf{tutoriales interactivos} y \textbf{consultas por voz}.
    \item \textbf{Accesibilidad}: avanzar al nivel AAA y hacer la revisión con lector de pantalla, que las pruebas automáticas no reemplazan.
    \item \textbf{Seguridad}: autenticación de dos factores para administradores; ampliar \texttt{audit\_log} a inicios de sesión, intentos fallidos, exportaciones y consultas al asistente; extender el cifrado, que hoy cubre las claves de API, a otros datos reservados.
    \item \textbf{Investigación}: repetir la prueba con usuario (Capítulo \ref{chap:pruebas}) en al menos diez PYMEs con la escala SUS; comparar modelos de lenguaje con la misma batería de \texttt{rag:evaluar}; medir el tiempo que ahorra el asistente; y documentar el método de desarrollo con asistentes de IA y sus controles.
\end{itemize}

\subsection{Hoja de ruta priorizada}

La Tabla \ref{tab:roadmap} presenta un orden sugerido. La prioridad combina el impacto esperado para una PYME como la del caso de estudio con la factibilidad, entendida como los recursos, las dependencias técnicas y el tiempo que requiere cada mejora.

\begin{table}[!htbp]
\centering
\caption{Hoja de ruta priorizada de trabajos futuros.}
\label{tab:roadmap}
\begin{tabular}{|c|>{\raggedright\arraybackslash}p{6cm}|c|c|}
\hline
\textbf{Prioridad} & \textbf{Trabajo Futuro} & \textbf{Impacto} & \textbf{Factibilidad} \\
\hline
1 & Copias de seguridad fuera del servidor & Alto & Alta \\
2 & Desglose del consumo por destino (RF-15) & Alto & Alta \\
3 & Pronóstico con series temporales y detección de anomalías & Medio & Media \\
4 & Pruebas de extremo a extremo y análisis estático en la integración continua & Medio & Alta \\
5 & Memoria conversacional en la recuperación & Medio & Alta \\
6 & Registro sin conexión en la bodega & Medio & Media \\
7 & Módulo de producción (BOM) & Medio & Media \\
8 & Clasificación de intenciones con NLP & Medio & Media \\
9 & Arquitectura multi-tenencia & Alto & Baja \\
10 & Integración con facturación electrónica DIAN & Alto & Baja \\
\hline
\end{tabular}
\end{table}
